\documentclass[11pt]{article}
\usepackage[T1]{fontenc}
\usepackage[utf8]{inputenc}
\usepackage{lmodern,amsmath,amssymb,amsthm}
\usepackage[margin=1in]{geometry}
\usepackage{hyperref}
\hypersetup{colorlinks=true,linkcolor=blue,urlcolor=blue}
\newtheorem{theorem}{Theorem}
\newtheorem{lemma}[theorem]{Lemma}
\newtheorem{proposition}[theorem]{Proposition}
\newtheorem{corollary}[theorem]{Corollary}
\theoremstyle{remark}
\newtheorem{remark}[theorem]{Remark}
\newcommand{\an}{\mathrm{an}}
\newcommand{\dd}{\mathrm d'\mathrm d''}
\newcommand{\Psh}{\operatorname{Psh}}
\newcommand{\Int}{\operatorname{Int}}
\newcommand{\Pic}{\operatorname{Pic}}
\newcommand{\trop}{\operatorname{trop}}
\title{P02.4: Ordinary smooth decreasing regularization\\
Verified sufficient curve sectors and an actual boundary obstruction}
\author{Fresh public-source argument; unverified mathematical draft}
\date{4 October 2026}
\begin{document}
\maketitle
\begin{abstract}
The original characterization is not solved. We give a stage-one ordinary
Chambert--Loir--Ducros (CLD) construction for every arbitrary classical
plurisubharmonic (psh) target on specified smooth curves over fields with a
countable dense subset, and extend the germ statement through finite unibranch
normalization. A quantitative shift makes independently smoothed lisse terms
decrease on one fixed ordinary neighborhood. We also treat all targets on
saturated products with proper projective fibres, without claiming this treats
all sheaf germs on those higher-dimensional spaces. At an actual analytic
boundary we construct a classical psh spike which cannot even be a local
pointwise limit of a sequence of continuous functions when the residue field
is uncountable. For this particular spike, countability of the residue field
is also sufficient, by an explicit ordinary smooth psh sequence. None of these
results gives the requested sharp characterization of all spaces.
\end{abstract}

\section{Scope, definitions, and checked published premises}
Only P02.4 is addressed. The required target is a decreasing \emph{sequence}
of finite ordinary CLD-smooth psh functions on one fixed ordinary open
neighborhood, converging pointwise to an arbitrary classical psh function
with values in $\mathbb R\cup\{-\infty\}$. Continuity, finiteness, or torus
invariance of the starting function is not imposed in the curve theorems.
The word classical has exactly the curve-test meaning in
\cite[Definition 4.1]{GR}: for every analytic field extension and every
analytic curve mapping into the extension of the space, the pullback is
subharmonic. The subharmonic definition allows the identically
$-\infty$ function; on a connected smooth strictly analytic boundaryless
curve it agrees with Thuillier's definition or the identically $-\infty$
case, by \cite[3.4.2 and 4.1.2]{GR}.

Ordinary smooth functions are represented on \emph{ordinary open} charts by
smooth functions of finitely many logarithms of invertible analytic functions.
This is the ordinary-open assertion of the actual saved CD v2,
\cite[Proposition 15.1.1(b), Corollary 15.1.2]{CD}. An ordinary smooth psh
function has $\dd f\geq0$ as a smooth $(1,1)$-form. In that saved v2 the
smooth stage is $\Psh_{X,0}$, the first decreasing-sequence stage is
$\Psh_{X,1}$, and these are defined in Lemma 8.1.4; Definition 8.1.13 calls
the latter locally approachable. Proposition 8.1.5 preserves stages under
pullback. It does not collapse stage two to stage one. On boundaryless
spaces these smooth functions also have positive $\dd[f]$ currents by
\cite[\S8.3, Proposition 8.3.1]{CD}; the boundaryless hypothesis in that
section is indispensable and is retained here.

The following exact published inputs were read, with their literal proofs
retained in the accompanying source extracts.
\begin{enumerate}
\item Thuillier's Theorem 3.4.2 \cite{Thu} gives a decreasing \emph{family/net}
of finite lisse subharmonic functions on each relatively compact ordinary
open subset of a smooth strictly analytic curve. Its proof uses finite
subsets of the non-rigid locus and harmonic interpolation. It does not give
the arbitrary-target sequence needed here.
\item Wanner's standing field in \cite{Wan} is complete, algebraically
closed, and nontrivially valued. Lemma 5.2 says: on a smooth proper algebraic
curve $C$ satisfying $H_C\subset C^\infty$, a lisse function on a strictly
affinoid $Y$ that is subharmonic on $Y\setminus\partial Y$ admits continuous
approximants on $Y$, ordinary smooth and psh on that same interior, converging
uniformly on $Y$. They can be chosen decreasing. Her Corollary 5.4 verifies
the harmonic-smoothness premise if the residue field is algebraic over a
finite field, or $C=\mathbb P^1$, or $C$ is a Mumford curve. The corollary's
target itself is \emph{finite continuous}; our arbitrary-target claim below
uses Lemma 5.2, not an unqualified invocation of that corollary.
\item GJR Proposition 15.11 \cite{GJR} has the conclusion ``a decreasing net
of weakly smooth psh functions'' on relatively compact opens of a rig-smooth
boundaryless curve over a nontrivially valued field. The proof first reduces
to one lisse term and constructs a sequence for that term. This does not
extract a sequence from the original net. Harmonic neighborhoods in that
proof provide weak, and not automatically ordinary, smoothness.
\item \cite[Propositions 3.5 and 4.2]{GR} gives sheafiness, field extension,
and analytic pullback of subharmonic and classical psh functions; finite
surjective maps also detect these properties. In particular Proposition
3.5(3) applies to a USC curve function subharmonic on the analytic interior,
including its arbitrary boundary values. The proof explicitly treats the
new interior components over a boundary point after scalar extension.
Theorem 6.1 gives the local maximum principle at analytic interior points.
The ordinary-smooth/classical comparison is the published argument in
\cite[1.2]{GR}, reducing curve tests to Wanner's continuous comparison on
smooth curve neighborhoods; it is not an unpublished book premise.
\end{enumerate}

Thuillier uses the all-domain site $X_G$ and the strictly analytic rigid
site $X_R$ separately (\cite[2.1.7]{Thu}); Definition 3.2.1 extends lisse
functions to $X_G$. His Proposition 3.2.4 gives the skeletal description
on compact \emph{strict} domains. GJR 15.8--15.9 likewise imposes strict
domains and the indicated vertex/value-group restrictions. We never infer
strict real-PL membership for all-domain lisse functions. We never replace
weakly smooth by ordinary smooth. No real-PL or Fubini--Study conclusion is
needed in this argument.

\section{Two elementary lemmas with quantitative bounds}
\begin{lemma}[Countable extraction]
Let $Z$ be second countable and let $(f_i)_{i\in I}$ be a decreasing net of
continuous real functions on $Z$ with pointwise infimum
$u:Z\to\mathbb R\cup\{-\infty\}$. There are indices $i_j$, increasing in
the directed ordering, such that $f_{i_j}$ decreases pointwise to $u$.
\end{lemma}
\begin{proof}
For every rational $q$, $\{u<q\}=\bigcup_i\{f_i<q\}$. A subspace of a
second-countable space is second countable and hence Lindelof: choosing one
covering open set for each basis element contained in a covering open set
gives a countable subcover. Thus choose countably many indices for each
$q$ covering $\{u<q\}$. List the union of all selected indices as $a_j$,
padding a finite list if necessary. Recursively take $i_j$ dominating
$i_{j-1}$ and $a_j$. Then $f_{i_j}\geq u$ and the sequence decreases.
If $u(z)<q$, some selected $f_{a_k}(z)<q$, so eventually
$f_{i_j}(z)<q$. All rational $q>u(z)$ show the limit is $u(z)$, and all
rational $q$ handle $u(z)=-\infty$. If no indices need selecting, the
conclusion is trivial. We did not claim the chosen indices are cofinal in
the entire directed set.
\end{proof}

\begin{lemma}[Strict affinoids over separable fields]
If $K$ has a countable dense subset, every strictly $K$-affinoid space is
second countable.
\end{lemma}
\begin{proof}
Choose a countable dense subring of $K$ by adjoining a countable dense
subset to its prime ring. The polynomials with these coefficients are a
countable dense subset of $K\langle T_1,\ldots,T_N\rangle$. Their images
give a countable dense subset $A_0$ of any strict affinoid quotient $A$,
with a submultiplicative quotient Banach norm. Every bounded multiplicative
seminorm $x$ on $A$ satisfies $|a(x)|\leq\|a\|$: apply boundedness to
$a^n$ and take $n$-th roots. The ultrametric inequality gives
$\bigl||a(x)|-|b(x)|\bigr|\leq\|a-b\|$ uniformly in $x$.
Every evaluation $|a|$ is consequently a uniform limit of evaluations
$|b|$ for $b\in A_0$. The topology on $\mathcal M(A)$ is therefore induced
by countably many real-valued evaluations. Rational intervals for those
coordinates and their finite intersections give a countable base.
\end{proof}

\begin{lemma}[One-stage smoothing with a fixed interstage gap]
Suppose $h_j:Y\to\mathbb R$ decreases pointwise to $u$, and for each
$j\geq1$ a finite smooth psh $s_j$ on the same $V\subset Y$ satisfies
$\|s_j-h_j\|_V\leq\epsilon_j$, where $\epsilon_j=2^{-j}$.
Then
\[
 v_j=s_j+4\epsilon_j
\]
is smooth psh on $V$, tends pointwise to $u|_V$, and
\begin{align*}
 h_j+3\epsilon_j&\leq v_j\leq h_j+5\epsilon_j,\\
 v_j-v_{j+1}&\geq(h_j-h_{j+1})+
 4(\epsilon_j-\epsilon_{j+1})-\epsilon_j-\epsilon_{j+1}
 \geq\epsilon_j/2>0.
\end{align*}
\end{lemma}
\begin{proof}
Adding a constant preserves smoothness and $\dd$-positivity.
The displayed bounds follow directly from the error inequalities and
$\epsilon_{j+1}=\epsilon_j/2$. At a finite target value the error tends to
zero, and at a $-\infty$ value the upper bound tends to $-\infty$.
All inequalities are pointwise on one fixed $V$. In particular this directly
produces stage-one approximants; no stage-two closure argument is used.
\end{proof}

\section{Arbitrary targets on a concrete curve sector}
\begin{theorem}\label{thm:curves}
Let $K$ be complete, algebraically closed and nontrivially valued, with a
countable dense subset. Let $C$ be a smooth proper algebraic curve satisfying
one of the following concrete sufficient conditions:
\begin{enumerate}
\item $\widetilde K$ is algebraic over a finite field;
\item $C=\mathbb P^1_K$;
\item $C$ is a Mumford curve.
\end{enumerate}
For every ordinary open $W\subset C^\an$, every arbitrary classical psh
$u:W\to\mathbb R\cup\{-\infty\}$, and every $x\in W$, there is a strictly
affinoid neighborhood $Y\subset W$ of $x$, with ordinary open interior
$V=Y\setminus\partial Y$, and a sequence of finite ordinary CLD-smooth
psh functions $v_j$ on $V$ decreasing pointwise to $u|_V$.
The construction is compatible with restriction, ordinary chart overlaps,
and pullback to every analytic extension field. If $u$ is finite continuous,
convergence is locally uniform; that extra conclusion is not asserted for
an arbitrary target.
\end{theorem}
\begin{proof}
Work in a connected ordinary neighborhood of $x$ in $W$. If the target is
identically $-\infty$ there, take $v_j=-j$. Otherwise it is subharmonic
in Thuillier's sense by \cite[3.4.2 and 4.1.2]{GR}. Choose a relatively
compact ordinary open $W'$ containing $x$ with closure in $W$.
Thuillier's Theorem 3.4.2 gives a decreasing net of finite lisse
subharmonic functions $f_i$ on $W'$ with infimum $u$.
Choose one strictly affinoid neighborhood $Y$ of $x$ contained in $W'$.
Such neighborhoods exist on these good strictly analytic algebraic curves;
their boundaries are finite and their interiors are ordinary open.

All $f_i|_Y$ are continuous. The second-countability and extraction lemmas
produce one decreasing sequence $h_j=f_{i_j}|_Y$ with infimum $u|_Y$.
This is the only countable extraction in the proof and its hypothesis has
been proved from the field condition. No continuity of $u$ is used.

For every $j$, $h_j$ is lisse on the strictly affinoid $Y$ and
subharmonic on the \emph{same} $V=Y\setminus\partial Y$.
The geometric conditions in the theorem give $H_C\subset C^\infty$ by
the checked proof of Wanner Corollary 5.4 and Thuillier Theorem 2.3.21.
Apply Wanner Lemma 5.2 to $h_j$ on $Y$. Select one of its continuous
approximants $s_j$ with
$\|s_j-h_j\|_Y\leq2^{-j}$. It is ordinary smooth and psh on $V$.
Wanner's construction uses harmonic pieces, which are ordinary smooth
under $H_C\subset C^\infty$, and convex regularized maxima on the connecting
annuli; its equalities with the original term near the gluing boundaries
establish ordinary-open gluing. Positivity is the positivity in Lemma 5.2,
not an inference from the word smooth.

Apply the fixed-gap lemma to obtain $v_j=s_j+4\cdot2^{-j}$. This is a
single decreasing sequence of ordinary smooth psh functions on $V$, with
gap at least $2^{-j-1}$ and pointwise limit $u$, including all permitted
$-\infty$ values and discontinuities. On a compact subset where the target
is finite continuous, Dini's theorem gives uniform convergence.

Each $v_j$ is a globally defined function on $V$. Its ordinary chart
representatives consequently agree as functions and as the induced
forms on intersections. We did not construct unrelated sequences on
different charts and then assume they agree. Restricting the sequence
to any ordinary open subset preserves all inequalities and its limit.
The field-extension and positivity details are supplied in the next
section and apply to these particular globally defined terms.
\end{proof}

The field hypothesis is a sufficient hypothesis, not a claimed necessary
one; the familiar field $\mathbb C_p$ is an example. The theorem treats all
classical psh sheaf germs on ordinary opens of the specified curves, rather
than only finite continuous global functions. Wanner's harmonic-smoothness
conditions have not been proved necessary for approximation.

\section{Ordinary pullbacks, scalar extension, and finite unibranch descent}
The pullback assertion here concerns ordinary smooth functions with positive
$\dd$ forms, and not arbitrary positive currents. If $g:Z\to V$ is analytic,
an ordinary open chart
\[
 v_j=\phi_j(-\log|a_1|,\ldots,-\log|a_r|),\qquad a_i\in\mathcal O^\times,
\]
pulls back to the same smooth real function of the invertible functions
$a_i\circ g$. Thus it is ordinary smooth. Differentials commute with
pullback, and pullback preserves smooth $(1,1)$-positivity by the linear
algebra of \cite[Lemma 1.2.11]{CD} and the ordinary tropical-chart
construction; equivalently the smooth case of CD Proposition 8.1.5.
This proves $\dd(g^*v_j)=g^*(\dd v_j)\geq0$. This is an application to
forms; we do not invoke a pullback operation for arbitrary currents.

For an analytic extension $L/K$ and $\pi:V_L\to V$, the charts pull back
to invertible $a_{i,L}$ on the ordinary opens $\pi^{-1}(U)$ and the same
$\phi_j$. To justify positivity without assuming that this particular
representing Hessian is positive on all of the ambient real space, apply
CD Lemma 6.5.1 to the positive ordinary form $\dd v_j$. It supplies an
ordinary neighborhood $U$, a tropical chart $f$, and a positive form
$\alpha$ on a paralin\'eaire set $P$ containing $f(U)$, with
$\dd v_j=f^*\alpha$. Refining the chart by Proposition 15.1.1(b)
expresses its coordinates as affine combinations of logarithms of
analytic units. Those coordinates therefore pull back under scalar
extension as well. The image of $f_L=f\circ\pi$ lies in the same $P$,
and the positive-linear-pullback test gives $f_L^*\alpha\geq0$.
Since differentials commute with this chart pullback,
$\dd(v_j\circ\pi)=f_L^*\alpha$. This proves positivity even when $L$
has new residue directions; it is not a harmonic-smoothness assertion
about all functions over $L$.
We obtain ordinary smooth psh $v_{j,L}=v_j\circ\pi$ and
\[
 v_{j,L}-v_{j+1,L}\geq2^{-j-1},\qquad
 \lim_jv_{j,L}(y)=u(\pi(y)).
\]
No second countability of $V_L$ is needed. The comparison in \cite[1.2]{GR}
makes these smooth psh functions classical; its field-compatible curve
test agrees with the positive-form pullback argument. On boundaryless
domains, CD \S8.3 also gives their positive currents. Successive analytic
extensions give exactly the same composed functions. Equality on overlaps
is preserved by taking preimages, so the scalar and overlap assertions
are actual equalities, not new choices of approximants.
This does \emph{not} prove that every new classical psh target over $L$
comes by descent from $K$, or that ordinary smoothness descends along an
arbitrary scalar extension.

\begin{theorem}[A singular unibranch curve sector]\label{thm:unibranch}
Let $K$ satisfy the field assumptions of Theorem \ref{thm:curves}.
Let $D$ be an integral proper algebraic curve whose finite normalization
$\nu:C\to D$ is bijective on geometric points and whose smooth proper
normalization $C$ satisfies one of that theorem's three conditions.
Then every arbitrary classical psh sheaf germ on $D^\an$ has a decreasing
ordinary CLD-smooth psh sequence on one fixed ordinary neighborhood.
The asserted sequences have the same restriction and scalar compatibility
as above. In particular this applies to cuspidal proper curves with
normalization $\mathbb P^1$.
\end{theorem}
\begin{proof}
On the smooth locus use Theorem \ref{thm:curves} through the normalization
isomorphism. The exceptional locus of a normalization of an integral
proper algebraic curve is a finite set of closed points; over this $K$
those are $K$-rational rigid points. Around a singular point $x$, remove
the other exceptional points. There is exactly one point $x'$ over $x$,
and $\nu$ is an isomorphism away from $x$. The analytic finite map is
closed and surjective. It is now bijective on the underlying neighborhood,
so it is a homeomorphism there. Indeed, outside $x$ it is an isomorphism,
and the fibre over $x$ has just the one analytic point $x'$.

Let $u$ be the target on an ordinary neighborhood $W$ of $x$.
Then $\widetilde u=u\circ\nu$ is classical psh by GR Proposition 4.2(6).
Construct the curve sequence $\widetilde v_j$ on an ordinary neighborhood
$V'$ of $x'$ by Theorem \ref{thm:curves}. Finiteness and closedness of
$\nu$ give an ordinary neighborhood $N\subset W$ of $x$ such that
$\nu^{-1}(N)\subset V'$: remove the closed image of the complement of
$V'$ in $\nu^{-1}(W)$. Descend the functions via the homeomorphism to
finite continuous functions $v_j$ on this single fixed $N$.

Every ordinary CLD-smooth function is locally constant at a $K$-rational
rigid point. To see this, use an ordinary chart with finitely many local
invertible functions $a_i$. On a sufficiently small neighborhood,
$|a_i-a_i(x')|<|a_i(x')|$ for every $i$, by analytic continuity and
$a_i(x')\ne0$. Hence all $|a_i|$ are constant there, so is the represented
smooth function. Therefore $\widetilde v_j$ is constant on some open
neighborhood of $x'$, whose image under the homeomorphism is open at $x$.
Its descended $v_j$ is constant on that neighborhood, and is ordinary
smooth elsewhere by the analytic isomorphism. This is ordinary-open
descent, not general descent of smooth forms along a finite map.

On the constant neighborhood $\dd v_j=0$; off $x$ its $\dd$ form is
nonnegative by the isomorphism with the normalization. These opens cover
$N$, so $v_j$ is smooth psh. For classical positivity one can also apply
the finite-surjective detection statement GR Proposition 4.2(6) to
$\nu^{-1}(N)\to N$. The functions agree on the intersections by their
global definition. Their inequalities, interstage gaps and limits are
those of $\widetilde v_j$, so $v_j\downarrow u|_N$. Their scalar pullbacks
are ordinary smooth psh and converge to $u_L|_{N_L}$ by the preceding
chart argument; this does not require a new assertion about normalization
commuting with every scalar extension.
\end{proof}

The one-point fibre is essential to this proof. With several branches,
separately chosen smooth sequences can have different constant values
near the preimages of the node at every finite stage. Agreement only of
their limits does not permit their descent. No claim is made for arbitrary
multibranch singularities, general finite maps, or strict scalar descent
over a non-algebraically-closed base field.

\section{All targets on a noninvariant saturated higher-dimensional sector}
\begin{proposition}\label{prop:saturated}
Under the smooth curve hypotheses of Theorem \ref{thm:curves}, let
$U\subset C^\an$ be ordinary open, $m\geq1$, and let
$X=U\times_K(\mathbb P^m_K)^\an$ with projection $p$.
For every arbitrary classical psh function $u$ defined on this entire
saturated $X$, and each $z\in U$, there is an ordinary neighborhood
$V$ of $z$ and a sequence of finite ordinary smooth psh functions on
the single ordinary open $p^{-1}(V)$ decreasing to $u|_{p^{-1}(V)}$.
There is no torus-invariance or continuity assumption on $u$.
\end{proposition}
\begin{proof}
For $z\in U$, the analytic fibre is
$(\mathbb P^m_{\mathcal H(z)})^\an$. Apply scalar extension and analytic
pullback, GR Proposition 4.2(5,6), to see that $u$ restricts to a classical
psh function on this fibre. Projective space analytification is compact,
connected and boundaryless. A USC function into $\mathbb R\cup\{-\infty\}$
on it is either identically $-\infty$ or has a finite attained maximum $M$:
local finite upper bounds give a global bound by compactness, and nested
closed upper-level sets give attainment. In the latter case its maximum
locus is nonempty and closed. At every point of that locus GR Theorem 6.1
applies, since the whole fibre is boundaryless, so the locus is open.
Connectedness implies that the function is constant. Thus $u$ is constant
on every proper fibre, including the entirely $-\infty$ fibres.

Choose a $K$-rational point $o\in\mathbb P^m(K)$ and its analytic section
$s:U\to X$. Set $w=u\circ s$. It is USC and classical psh by analytic
pullback, hence subharmonic on the curve $U$. Fibrewise constancy gives
the equality of functions $u=w\circ p$, including infinite values.
Apply Theorem \ref{thm:curves} to $w$ at $z$ on a fixed $V$. Its smooth
psh terms $w_j$ pull back to ordinary smooth psh $w_j\circ p$ by the
ordinary pullback proof above, and
\[
 w_j\circ p-w_{j+1}\circ p\geq2^{-j-1},\qquad
 w_j\circ p\downarrow w\circ p=u.
\]
The identically $-\infty$ case uses $-j$ instead. The entire inverse
image $p^{-1}(V)$ is one ordinary open neighborhood of every point over
$z$. The global formula establishes overlap agreement. After every
analytic extension it remains the same formula with $p_L$ and the
pulled-back $w_j$; its positivity and decreasing convergence remain valid.
The products here are boundaryless, so the smooth current is positive as
well.
\end{proof}

This is an assertion about \emph{all targets on saturated domains}.
It is not the assertion that every sheaf germ on $X$ has been regularized.
A target initially defined only on a nonsaturated affine chart cannot
be restricted to a whole proper fibre, and the preceding maximum-principle
argument is unavailable. No extension of that target to the saturated
domain has been proved. Consequently this proposition does not establish
the original P02.4 local sheaf property for these higher-dimensional
spaces. It does give a noninvariant higher-dimensional sufficient sector
with arbitrary, possibly discontinuous or $-\infty$-valued, targets in
the specified domain of definition.

\section{A sharp sequence obstruction at an actual analytic boundary}
Let $K$ now be any complete algebraically closed nontrivially valued
field; no dense-countability assumption is imposed. Write
$D=\mathcal M(K\langle T\rangle)$, let $\eta$ be its Gauss point, and
fix $A>0$. Define
\[
 u(\eta)=A,\qquad u(x)=0\quad(x\ne\eta).
\]
The singleton $\{\eta\}$ is closed in the Hausdorff disc, so $u$ is USC.
The analytic boundary of this closed disc is exactly $\{\eta\}$ and
$u=0$ on $\Int(D)$. GR Proposition 3.5(3), with its scalar-compatible
boundary proof, says $u$ is subharmonic on $D$. GR 4.1.2 makes it
classically psh. This step verifies the actual starting class, rather
than just exhibiting a topological USC function.

For each $\alpha\in\widetilde K$, choose a representative
$a_\alpha\in K^\circ$ and set
\[
 D_\alpha=\{x\in D:|T-a_\alpha|(x)<1\}.
\]
These are nonempty pairwise disjoint ordinary open residue discs, and
partition $D\setminus\{\eta\}$. One elementary verification of the last
assertion is as follows. If every $|T-a_\alpha|(x)=1$, every nonzero
polynomial, factored over the algebraically closed $K$, has its Gauss
norm at $x$: roots in $K^\circ$ give factors of norm one and roots
outside $K^\circ$ give their constant dominating absolute value.
Approximate Tate series uniformly by polynomials to conclude that $x$
is the Gauss seminorm. Conversely $|T-a_\alpha|(x)<1$ singles out exactly
that residue class by the ultrametric inequality.

Every ordinary neighborhood of $\eta$ contains all but finitely many
\emph{whole} $D_\alpha$. Indeed a basic neighborhood uses finitely many
evaluations of Tate series $f$. Normalize a nonzero $f$ by its Gauss norm;
its reduction is a nonzero polynomial over $\widetilde K$, since its
coefficients tend to zero. Outside the finitely many roots of that
reduction, $|f|$ is its Gauss norm throughout the corresponding residue
disc. For a direct verification, translate by $a_\alpha$: if its reduced
value there is nonzero, the constant term has full norm, while all
nonconstant terms times $(T-a_\alpha)^n$ have strictly smaller norm
at every point of that open disc. Thus any finite collection of such
evaluations agrees with its value at $\eta$ outside finitely many
directions. This proves the neighborhood assertion.

\begin{theorem}[Uncountable residual boundary obstruction]\label{thm:obstruction}
If $\widetilde K$ is uncountable, the above classical psh $u$ is not a
pointwise limit of \emph{any sequence of finite continuous functions}
on any ordinary neighborhood of $\eta$ in $D$. In particular $D$ fails
the original decreasing ordinary CLD-smooth psh property at $\eta$.
\end{theorem}
\begin{proof}
Suppose continuous $v_j:N\to\mathbb R$ converge pointwise to $u|_N$,
where $N$ is an ordinary neighborhood of $\eta$. Discard finitely many
indices so that $v_j(\eta)>3A/4$. For each remaining $j$, continuity
gives an ordinary neighborhood $N_j\subset N$ of $\eta$ on which
$v_j>A/2$. By the preceding neighborhood assertion, there is a finite
set $E_j\subset\widetilde K$ such that $D_\alpha\subset N_j$ for
$\alpha\notin E_j$. Choose
$\alpha\notin\bigcup_jE_j$, possible by uncountability, and its rigid
point $a_\alpha$. All $v_j(a_\alpha)>A/2$, but $u(a_\alpha)=0$, a
contradiction. Decreasingness and psh positivity were not even needed.
\end{proof}

\begin{corollary}[Higher-dimensional failures]
For the same $K$ with uncountable residue field and any $m\geq1$,
$D\times_K(\mathbb P^m_K)^\an$ fails P02.4 at
$(\eta,o)$ for any $K$-rational $o$.
\end{corollary}
\begin{proof}
The pullback $p^*u$ under projection is classical psh by GR Proposition
4.2(6). If continuous, in particular ordinary smooth, approximants
existed on a neighborhood $N$ of $(\eta,o)$, restrict them to the
analytic section $D\times\{o\}$. Its inverse image of $N$ is an
ordinary neighborhood of $\eta$, and its target is exactly $u$.
Theorem \ref{thm:obstruction} forbids even a continuous sequence there.
These are good strictly analytic spaces with a genuine analytic boundary;
the obstruction is not a purported residue-Picard obstruction on a
boundaryless curve.
\end{proof}

\begin{proposition}[Sharpness for this spike]\label{prop:countable-spike}
If $\widetilde K$ is countable, the same spike admits a decreasing
sequence of finite ordinary CLD-smooth functions with nonnegative $\dd$
form on the \emph{entire fixed closed disc} $D$. Thus, for this target,
countability of the residue field is both necessary and sufficient.
These are the ordinary smooth stage-zero psh functions of CD.
\end{proposition}
\begin{proof}
Enumerate $\widetilde K$ and choose representatives $a_1,a_2,\ldots$.
(The residue field of an algebraically closed field is infinite, so no
finite-list issue arises.) Choose a nonzero nonnegative smooth bump
$\beta$ supported in $(-1,-1/2)$. Set
\[
 c=\frac{A}{\int_{-1}^{-1/2}(-t)\beta(t)\,dt},\qquad
 \chi(s)=c\int_{-\infty}^s(s-t)\beta(t)\,dt.
\]
Then $\chi$ is finite smooth convex nondecreasing, $\chi(s)=0$ for
$s\leq-1$, $\chi(0)=A$, and $0\leq\chi(s)\leq A$ for $s\leq0$.
Put $P_j(T)=\prod_{i=1}^j(T-a_i)$ and
\[
 v_j(x)=\chi\bigl(j\log|P_j(x)|\bigr),
 \quad\text{with }v_j(x)=0\text{ if }P_j(x)=0.
\]
On the ordinary open locus $P_j\ne0$, this is a smooth function of the
logarithm of an invertible analytic function. Near a zero, the open
set $|P_j|<e^{-1/j}$ has $v_j=0$, so this extension is ordinary smooth
there as well. These ordinary open sets cover $D$ and their functions
agree on intersections. On the nonzero locus, write $\ell=j\log|P_j|$;
the logarithm of a unit is affine in an ordinary tropical chart, so
\[
 \dd v_j=\chi''(\ell)\,\mathrm d'\ell\wedge\mathrm d''\ell\geq0.
\]
It vanishes on the constant neighborhoods of zeros. Hence the functions
are smooth psh in the CD stage-zero positive-form sense. Positive currents
are also obtained on their analytic interiors, but CD \S8.3 is not cited
to assert an unqualified current statement at the boundary.

On $D$ every $\log|T-a_i|\leq0$. On the common nonzero locus,
\[
 (j+1)\log|P_{j+1}|-j\log|P_j|
 =\sum_{i\leq j}\log|T-a_i|+(j+1)\log|T-a_{j+1}|\leq0.
\]
Monotonicity of $\chi$, together with the zero values when a factor
vanishes, proves $v_{j+1}\leq v_j$. At $\eta$ every factor has norm
one, so $v_j(\eta)=A$. If $x\ne\eta$, choose $i$ with
$|T-a_i|(x)=r<1$. For $j\geq i$ with no vanishing factor,
$j\log|P_j(x)|\leq j\log r\to-\infty$. Thus $v_j(x)$ is eventually
exactly zero; a vanishing factor gives the same conclusion immediately.
Therefore $v_j\downarrow u$ on this one fixed domain.

After an arbitrary analytic extension $L/K$, the identical polynomials,
cutoff neighborhoods and positive Hessian formula give ordinary smooth
psh $v_{j,L}$. Pointwise composition gives
$v_{j,L}\downarrow u\circ\pi_{L/K}$. If $\widetilde L$ is uncountable,
this is not a contradiction: the target is the spike on the full
inverse image $\pi_{L/K}^{-1}(\eta)$, not the singleton-Gauss spike
on $D_L$. For example new residue directions over $L$ can map entirely
to the old Gauss point. The two targets must not be identified.
\end{proof}

The ``if and only if'' just proved concerns this one spike, not all
classical psh functions on all discs with countable residue field. In
particular it does not supply a sharp characterization for the original
arbitrary-target problem at analytic boundaries.

\section{Residue Picard/genus effects do not prove a failure of approximation}
In the precise semistable type-two setup of Thuillier Theorem 2.3.21 and
Lemma 2.3.22, the latter identifies
\[
 H_{C,x}/F_{C,x}\simeq\Pic^0(C_x)\otimes_{\mathbb Z}\mathbb R,
\]
where $C_x$ is the proper geometrically connected residual component and
$F$ is generated by logarithms of local analytic units. GJR Remark 15.10
identifies $F$ with the ordinary smooth harmonic functions. Thus a nonzero
quotient implies some harmonic germ is not ordinary smooth. It does not
assert that the germ cannot be a decreasing limit of ordinary smooth
\emph{psh}, nonharmonic functions. Wanner Remark 5.5 explicitly leaves
ordinary approximation on such curves unknown in that paper.

If the residual component is rational, $\Pic^0=0$. If the algebraically
closed residue field is algebraic over a finite field, every Jacobian
point is defined over some finite field and belongs to its finite group
of rational points, hence is torsion. Thus its tensor with $\mathbb R$
is zero even when the residual genus is positive. Positive residual genus
alone is not an obstruction. The actual PDF at printed page 59 says
``est de cardinal fini'' in the algebraic-finite-field bullet; we do not
use that stronger sentence about the entire group, which need not be
finite. The elementary pointwise torsion argument above is the needed
justification. Direct PDF inspection confirmed the Picard quotient,
the theorem's two conditions, and that GJR 15.11 really says a net.

We have proved neither necessity of $H_C\subset C^\infty$ for
approximation nor removal of it. In particular the genuine boundary
counterexample in the preceding section must not be reinterpreted as
a residue-genus counterexample on smooth boundaryless curves.

\section{Remaining exact original obligations and reproducible evidence}
The exact original P02.4 remains open in this draft. The results here do
not characterize higher-dimensional good strict analytic spaces for all
classical psh sheaf germs. Unresolved obligations include:
\begin{enumerate}
\item arbitrary noninvariant targets on nonsaturated higher-dimensional
ordinary charts and a geometric criterion covering all such targets;
\item a sharp positive/failure range beyond the displayed curve sector
and the displayed residual-cardinality boundary spike;
\item boundaryless curves with $H/F\ne0$, without turning one nonsmooth
harmonic germ into an unproved approximation obstruction;
\item arbitrary multibranch or nonalgebraic singular curves, and
ordinary smooth descent and synchronization along general finite normalization;
\item nonseparable fields in the arbitrary-target positive curve theorem,
where its explicit net-to-sequence extraction no longer applies;
\item non-algebraically-closed bases and strict scalar descent. Existing
scalar compatibility here concerns pullbacks of the constructed sequence,
not descent of independently constructed sequences;
\item compatibility needed to glue approximants from unrelated local
tropical charts. Here all asserted overlaps come from actual globally
defined terms on one neighborhood, not from a general gluing theorem.
\end{enumerate}

This is a mathematical argument based on checked published premises,
not an exhaustive finite search or computational proof. Executable
checks only validate retained inputs, source bytes, arithmetic of the
explicit interstage estimates, and document generation. They do not
establish mathematical correctness of the new arguments.

All files are in the fresh working directory. The original bytes are
retained in \texttt{input.json} and \texttt{input-original.json}; the
assignment and task are retained in \texttt{assignment-original.txt}
and \texttt{task-original.json}. The full proof is in
\texttt{P02.4-argument.md} and embedded in this standalone
\texttt{P02.4-paper.tex}. The seven authorized source PDF/extract pairs,
their supplied and actual hashes, selected literal statements and complete
proofs, exact application checks, and direct PDF inspection records are
listed in \texttt{source-extracts/source-provenance.json}.
There was no reading of earlier campaign workspaces or unpublished book
sources, no network use, no installation, no subagent or other-model
invocation, and no runtime changes or heavy computation.

Rerun from this directory with
\begin{verbatim}
python3 verify_sources.py
python3 build_and_check.py
\end{verbatim}
The second script invokes the existing \texttt{pandoc} to render the same
standalone proof into Markdown and the existing \texttt{pdflatex} twice,
with \texttt{-interaction=nonstopmode}, \texttt{-halt-on-error}, and
\texttt{-file-line-error}.
Actual return codes, warnings, error lines and complete transcripts are
retained in \texttt{checks.json}, \nolinkurl{build-diagnostics.json}, and
\texttt{build/}. The environment versions are retained in
\texttt{environment.json}. The PDF is a reading artifact; successful
builds and matching hashes do not certify any theorem.

\begin{thebibliography}{9}
\bibitem{GR} W. Gubler and J. Rabinoff,
\emph{Classically psh and pluriharmonic functions on Berkovich spaces},
saved public arXiv:2506.13548v2, especially 3.4--3.5, 4.1--4.2, 6.1.
\url{https://arxiv.org/pdf/2506.13548v2}.
\bibitem{GRPL} W. Gubler and J. Rabinoff, public arXiv:2503.05513,
saved supplied version; PL and piecewise smooth foundations checked,
not used as arbitrary ordinary smooth regularization.
\url{https://arxiv.org/pdf/2503.05513}.
\bibitem{Thu} A. Thuillier, thesis, 2005, saved tel-00010990v1,
especially 2.1.7, 2.3.21--2.3.22, 3.2.1--3.2.4, 3.4.2.
\url{https://theses.hal.science/tel-00010990v1/document}.
\bibitem{GJR} W. Gubler, P. Jell and J. Rabinoff,
\emph{Forms on Berkovich spaces based on harmonic tropicalizations},
saved public arXiv:2111.05741, 15.8--15.11.
\url{https://arxiv.org/pdf/2111.05741}.
\bibitem{Ber} V. G. Berkovich, \emph{Etale cohomology for non-Archimedean
analytic spaces}, Publ. Math. IHES 78 (1993), 5--161, supplied public cache.
\url{https://www.numdam.org/item/PMIHES_1993__78__5_0/}.
\bibitem{CD} A. Chambert-Loir and A. Ducros,
\emph{Formes differentielles reelles et courants sur les espaces de
Berkovich}, actual saved public arXiv:1204.6277v2. Numbering here refers
to that supplied v2, not the older numbering used in Wanner.
\url{https://arxiv.org/pdf/1204.6277v2}.
\bibitem{Wan} V. Wanner, \emph{Comparison of two notions of subharmonicity
on non-Archimedean curves}, saved public arXiv:1801.04713v3;
standing terminology, Lemma 5.2, Corollary 5.4, Remark 5.5.
\url{https://arxiv.org/pdf/1801.04713v3}.
\end{thebibliography}
\end{document}
